% =====================================================================================
% §2 METHODOLOGY — the full econometric treatment (classical GVAR and MFFM in parallel).
% Every constant here is verified against the code: cgl_bvarV5.m (hyperpriors, MH bounds),
% export_bvar_gvar24_loop.m (MAXIT/TOL/draws/WN columns/star construction), the shipped
% coupling artifact (dim, rho). Do not edit numbers by hand.
% =====================================================================================

Two models organize this section, and they are meant to be read against each other. The
classical GVAR comes first, in the form the literature developed it, because the MFFM keeps
its skeleton and the comparison is sharpest with both on the page. The MFFM follows: what a
bloc is, how its foreign variables are built, how the system is estimated and then solved.
The section closes with the question any design must answer --- what this way of working
buys, and what it costs --- and with the US bloc, which runs through everything as the
concrete example. Readers who want the tour first can jump to Section~\ref{sec:example} and
return.

\subsection{The classical GVAR}\label{sec:classical}

The global VAR grew out of a practical need: to extend a national model with the foreign
variables that national outcomes depend on, without attempting an intractable world system.
Proposed by \citet{PesaranSchuermannWeiner2004} and developed fully in
\citet{DeesDiMauroPesaranSmith2007}, it models each economy as a VARX$^{*}$, a VAR augmented
with weakly exogenous foreign variables, and ties the country models into one global system
through trade weights; \citet{ChudikPesaran2016} survey the approach and
\citet{diMauroPesaran2013} collect its applications.

Formally, country $i=1,\dots,N$ is described by a VARX$^{*}(p_i,q_i)$ model,
\begin{equation}\label{eq:varx}
  y_{it} \;=\; a_{i0} \;+\; \sum_{s=1}^{p_i} \Phi_{is}\, y_{i,t-s}
  \;+\; \sum_{s=0}^{q_i} \Lambda_{is}\, y^{*}_{i,t-s}
  \;+\; \sum_{s=0}^{q_i} \Upsilon_{is}\, d_{t-s} \;+\; u_{it},
\end{equation}
where $y_{it}$ is the $k_i\times1$ vector of domestic variables, $y^{*}_{it}$ the
country-specific foreign variables, $d_t$ a small set of global variables (the oil price is
the canonical member), and $u_{it}$ the reduced-form innovation. The foreign variables are
trade-weighted cross sections of the partners' variables,
\begin{equation}\label{eq:starsclassical}
  y^{*}_{it} \;=\; \sum_{j=1}^{N} w_{ij}\, y_{jt}, \qquad w_{ii}=0,
\end{equation}
with fixed weights $w_{ij}$, usually bilateral trade shares.\footnote{The weighting scheme
matters less than one might expect; see \citet[p.~169]{diMauroPesaran2013}.}

Two assumptions make this system estimable country by country. The first is \emph{weak
exogeneity} of $y^{*}_{it}$ and $d_t$ with respect to the parameters of bloc
$i$:\footnote{Foreign variables may affect the domestic economy contemporaneously, but they
are not affected by domestic disequilibria --- the lagged domestic error-correction terms do
not enter the marginal model of the foreign variables
\citep[pp.~56--57]{GarrattLeePesaranShin2012}. Careful applications test this restriction
bloc by bloc rather than assume it.} each economy takes the world as given. That is
plausible for a small open economy and fails, by construction, for the largest one, so the
United States is treated asymmetrically: it is the source of the global variables $d_t$,
and its own foreign-variable set is restricted. The second assumption is a common
estimation window. The blocs are estimated on a shared balanced sample, so ragged country
coverage is resolved by truncating everyone to the quarters all blocs share.

In its full form the classical model is a \emph{cointegrating} system. Each country model
is cast as a VECMX$^{*}$, the number of cointegrating relations is determined by
Johansen-type rank tests, estimation is by reduced-rank regression conditional on the
weakly exogenous variables, and over-identifying long-run restrictions can be imposed so
that the estimated equilibria match theory \citep{GarrattLeePesaranShin2012}. The short-run
coefficients are then estimated by least squares conditional on the long-run relations.
This long-run structure is one of the GVAR's real strengths: it disciplines the medium run
exactly where an unrestricted VAR is weakest.

The global solution comes from stacking. Collect $z_{it}=(y_{it}',\,y^{*\prime}_{it})'$ and
define the link matrices $W_i$, built from the weights, that extract bloc $i$'s own and
foreign variables from the global vector $y_t=(y_{1t}',\dots,y_{Nt}')'$, so that
$z_{it}=W_i\,y_t$. Writing \eqref{eq:varx} at first order for exposition and substituting,
\begin{equation}
  A_i W_i\, y_t \;=\; a_{i0} + B_i W_i\, y_{t-1} + u_{it},
  \qquad A_i=(I_{k_i},\,-\Lambda_{i0}), \quad B_i=(\Phi_{i1},\,\Lambda_{i1}),
\end{equation}
and stacking the $N$ blocs,
\begin{equation}\label{eq:stack}
  G\,y_t \;=\; a_0 + H\,y_{t-1} + u_t
  \qquad\Longrightarrow\qquad
  y_t \;=\; G^{-1}a_0 + G^{-1}H\,y_{t-1} + G^{-1}u_t,
\end{equation}
where $G$ stacks the $A_iW_i$ and $H$ the $B_iW_i$. The contemporaneous matrix $G$ is
invertible in applications, and inverting it \emph{once} delivers the world reduced form,
from which forecasts and generalized impulse responses \citep{PesaranShin1998} follow.
Stability requires the eigenvalues of $G^{-1}H$ to lie in the unit disc, a property that is
checked rather than guaranteed; applications routinely track how many eigenvalue moduli
approach one, since near-unit roots signal the over-parametrization risks inherent in
systems this large. Estimation in practice is supported by a mature toolchain
\citep{SmithGalesi2014}.

Figure~\ref{fig:pesaran} draws the pipeline: separate country estimations, stars from
observed data, one stack, one inversion. That is the machine the MFFM descends from.
Everything that follows keeps its two central ideas --- sparse trade-weighted linkage and a
stacked linear solution --- and changes how the blocs are specified, estimated, and solved.

\begin{figure}[t]
\centering
\includegraphics[width=0.92\textwidth]{figs/tikz/fig_pesaran.pdf}
\caption{The classical GVAR pipeline. Country models are estimated one at a time on their own
samples, linked by weakly exogenous trade-weighted stars built from observed data, then
stacked through the link matrices and inverted once to give the world reduced form
\eqref{eq:stack}. The two red tags mark the assumptions the MFFM replaces: weak exogeneity
and the common balanced estimation window.}
\label{fig:pesaran}
\end{figure}

\subsection{The MFFM country model}\label{sec:gbvar}

The MFFM strips the classical bloc down to a closed VAR and moves the difficulty elsewhere
--- out of the specification and into the estimation. Each of the $N=\NBlocs{}$ coupled
blocs is a closed, constant-coefficient VAR($p$),
\begin{equation}\label{eq:mffmbloc}
  y_{it} \;=\; c_i \;+\; \Gamma_i\, D_t \;+\; \sum_{\ell=1}^{p} B_{i\ell}\, y_{i,t-\ell}
  \;+\; u_{it},
  \qquad u_{it}\sim\mathcal N(0,\Sigma_i), \qquad p=\NLags,
\end{equation}
where $D_t=\bigl(\mathbf 1\{q(t){=}2\},\,\mathbf 1\{q(t){=}3\},\,\mathbf 1\{q(t){=}4\}\bigr)'$
holds quarterly seasonal dummies, with $\Gamma_i\neq0$ for the \NSeasonal{} blocs whose
source series are published without seasonal adjustment and $\Gamma_i=0$ elsewhere. The
bloc vector carries the partition
\begin{equation}\label{eq:partition}
  y_{it} \;=\; \bigl(\, x_{it}',\;\; g_t',\;\; x^{*\prime}_{it} \,\bigr)',
\end{equation}
with $x_{it}$ the domestic block (activity, prices, interest rates, credit, housing and,
for \NFiscal{} blocs, four fiscal drivers), $g_t$ the \NGlobals{} global financial prices
(the WTI oil price, the US ten-year Treasury yield, the US Baa corporate yield, and a US
equity index), and $x^{*}_{it}=(\mathrm{ROWDEM}_{it},\mathrm{ROWINF}_{it},\mathrm{ROWFIN}_{it})'$
the three foreign variables. Variables tagged as logarithmic enter as $100\ln$ levels, so
the displayed quarter-on-quarter annualized growth is $4\Delta$; rates and GDP ratios enter
untransformed. In the US bloc the four global prices are ordinary domestic columns --- the
United States is their source; every other bloc carries the same four series and receives
their paths from the US solution.

The exogenous side is deliberately leaner than the Bayesian GVAR it is closest to.
\citet{CrespoFeldkircherHuber2016} give each country a near-full foreign mirror,
$x^{*}_{it}=(y^{*},\allowbreak\Delta p^{*},\allowbreak eq^{*},\allowbreak i_s^{*},\allowbreak i_l^{*},\allowbreak tc^{*})'$ --- trade-weighted counterparts of
six of their seven domestic variables, all but the real exchange rate --- together with the oil
price as a global control determined inside the US model. The MFFM keeps only the three stars for
partner demand, inflation and the short rate, and carries the financial linkage instead in the
\NGlobals{} \emph{global} prices $g_t$: oil and three US-originated series (the ten-year yield, the
Baa corporate yield and an equity index) that enter every bloc directly rather than as
trade-weighted aggregates. That is the substantive difference between the two spillover mechanisms.
Where \citeauthor{CrespoFeldkircherHuber2016} transmit foreign equity, long rates and credit through
the weight matrix, the MFFM transmits US financial conditions as common prices, so a repricing at
the source reaches every bloc's own equations at once --- the mechanism behind the ``US credit event
as world credit event'' reading of Section~\ref{sec:example}. The leaner star block trades the
country-pair richness of a full foreign mirror for a sparser, more identifiable linkage and that
common-price channel.

The defining difference from \eqref{eq:varx} lies in what is \emph{absent}: there is no
contemporaneous term $\Lambda_{i0}x^{*}_{it}$ and no weak-exogeneity restriction. The
foreign variables are full endogenous columns with equations of their own, so each bloc is
a closed VAR that can forecast its foreign environment unaided. Contemporaneous
co-movement between foreign and domestic variables is carried by the covariance $\Sigma_i$
rather than by a structural loading, and the cross-country consistency that weak
exogeneity buys the classical model is imposed twice elsewhere: on the star \emph{data} at
estimation time (Section~\ref{sec:estimation}) and on the star \emph{paths} at solve time
(Section~\ref{sec:solve}). One consequence deserves emphasis. Nothing in this design stops
the United States from carrying foreign variables of its own, and it does: the classical
asymmetry disappears, and the world feeds back into the US bloc as into any other.

The star data are built as in \eqref{eq:starsclassical}, with one refinement forced by the
unbalanced panel. Let $\ell^{g}_{jt}$ and $\ell^{\pi}_{jt}$ denote partner $j$'s $100\ln$
real GDP and consumer-price levels and $r_{jt}$ its short rate (the policy rate where one
exists, the ten-year yield as fallback), all read from $j$'s \emph{completed} panel in the
sense of Section~\ref{sec:estimation}. The three foreign variables of bloc $i$ are
\begin{align}
  \mathrm{ROWDEM}_{it} &= 4\sum_{j\neq i}\tilde w^{\mathrm{DEM}}_{ij,t}
      \bigl(\ell^{g}_{j,t}-\ell^{g}_{j,t-1}\bigr), \label{eq:rowdem}\\
  \mathrm{ROWINF}_{it} &= 4\sum_{j\neq i}\tilde w^{\mathrm{INF}}_{ij,t}
      \bigl(\ell^{\pi}_{j,t}-\ell^{\pi}_{j,t-1}\bigr), \label{eq:rowinf}\\
  \mathrm{ROWFIN}_{it} &= \phantom{4}\sum_{j\neq i}\tilde w^{\mathrm{FIN}}_{ij,t}\; r_{j,t},
      \label{eq:rowfin}
\end{align}
the first two trade-weighted annualized growth rates, the third a level. The applied
weights renormalize the fixed bilateral goods-trade shares $w_{ij}$ (IMF merchandise-trade
statistics, 2021--23 average) over the partners available and eligible at each date, per
component:
\begin{equation}\label{eq:wtilde}
  \tilde w^{\,k}_{ij,t} \;=\;
  \frac{w_{ij}\,\mathbf 1\{j\ \text{eligible for}\ k\ \text{at}\ t\}}
       {\sum_{j'} w_{ij'}\,\mathbf 1\{j'\ \text{eligible for}\ k\ \text{at}\ t\}},
  \qquad k\in\{\mathrm{DEM},\mathrm{INF},\mathrm{FIN}\},
\end{equation}
so a partner whose sample has not yet started lends its weight to the others instead of
truncating everyone to a common window.\footnote{Eligibility differs by component: a bloc
with no usable short-rate series drops out of $\mathrm{FIN}$ only, and Argentina's 2023--24
inflation episode, which would dominate the price and rate averages, excludes it from
$\mathrm{INF}$ and $\mathrm{FIN}$ while it remains in $\mathrm{DEM}$.}

One modeling choice remains, and it is the deepest difference between the two designs: the
long run. The classical GVAR disciplines it with cointegrating restrictions; we discipline
it with shrinkage. The MFFM works in levels under a Minnesota-type prior
centered on univariate random walks, the Bayesian tradition for possibly integrated data
that avoids pre-testing for ranks \citep{DoanLittermanSims1984,Litterman1986,simszha1998}.
This shrinkage-in-levels route to the GVAR is itself established:
\citet{CrespoFeldkircherHuber2016} first replaced the classical cointegration step with
Bayesian shrinkage on the country-model coefficients, and showed that hierarchical
Minnesota-type and stochastic-search priors forecast an international macro-financial panel
more accurately than a random walk with drift, an unshrunk global model, or isolated country
VARs. The MFFM shares that estimation philosophy --- exogenous trade-weight links, a
VARX$^{*}$ per bloc, shrinkage in place of rank tests --- and differs in what the estimated
system is then put to: an interactive conditional solve (Section~\ref{sec:solve}), a joint EM
completion of the pandemic-holed panel, an endogenous anchor economy, and the fiscal and
external accounting satellites, rather than density forecasting alone.
The deliberate exceptions are the fiscal ratios --- revenue, primary expenditure, the
effective interest rate, the stock-flow term --- whose prior is centered on white noise
(own-lag coefficient zero): a GDP share is not a random walk, and a fiscal drift the data
do not insist on should not be allowed to compound over a five-year projection. What is
given up relative to the VECMX$^{*}$ is the explicit long-run structure; what is gained is
that no rank and restriction decisions have to be made, and hold, for
\NBlocs{} blocs simultaneously. For blocs with $\Gamma_i\neq0$ one further convention
matters in daily use: the deterministic term is quarter-specific in history but averaged in
projection --- forecasts replace $c_i+\Gamma_iD_t$ by
$\bar c_i = c_i + \tfrac{1}{4}(\gamma_{i,Q2}+\gamma_{i,Q3}+\gamma_{i,Q4})$ --- so projected
paths are seasonally adjusted while the model fits the unadjusted data it was given. A
related accrual convention spreads coupon-lumpy government interest payments evenly within
each calendar year for \NEvenEffi{} fiscal blocs; the online manual documents both.

\subsection{Estimation}\label{sec:estimation}

What makes estimation nonstandard here is not the prior but the data. The panel is
ragged: bloc samples differ in length by more than a decade, and at any moment
the latest quarter is published for some blocs and not others. The pandemic quarters
2020Q2--Q4 are treated as missing by choice, for the macroeconomic block only, so that
three observations of unprecedented size do not dominate the estimated dynamics of thirty
normal years \citep{LenzaPrimiceri2022}; the financial columns stay observed, and of the
foreign variables $\mathrm{ROWDEM}$ and $\mathrm{ROWINF}$ are masked with the macro block
while $\mathrm{ROWFIN}$ remains observed. And the regressors $x^{*}_{it}$ are not data at
all: they are functions of the partners' incomplete panels. Estimation therefore solves
three problems jointly within each bloc --- missing cells, hyperparameters, coefficients
--- and a fourth across blocs, because the stars each bloc needs are outputs of the other
blocs' completions.

Within a bloc, everything runs through the state-space form. For given coefficients, stack
$\alpha_t=(y_{it}',\dots,y_{i,t-p+1}')'$ and write
\begin{equation}\label{eq:ss}
  \alpha_t = \mathbf T_i\,\alpha_{t-1} + c_\alpha + \eta_t, \qquad
  y^{\mathrm{obs}}_{it} = Z_{it}\,\alpha_t + \varepsilon_{it}, \qquad
  \varepsilon_{it}\sim\mathcal N(0,\,10^{-12}I),
\end{equation}
where $\mathbf T_i$ is the companion matrix of \eqref{eq:mffmbloc} and the selection
$Z_{it}$ keeps exactly the cells observed at date $t$, so missing cells simply drop out of
the measurement equation. The simulation smoother of \citet{DurbinKoopman2002} then draws
the missing cells exactly, conditional on everything observed --- data augmentation in the
sense of \citet{TannerWong1987}.

The prior on $(\beta_i,\Sigma_i)$, where $\beta_i$ stacks
$(c_i,\Gamma_i,B_{i1},\dots,B_{ip})$, is the natural-conjugate normal--inverse-Wishart
family with the hierarchical Minnesota structure of \citet{Giannone2015}. Conditional on
$\Sigma_i$, the prior mean of each equation is the univariate random walk --- own first-lag
coefficient one --- except on the fiscal-ratio columns, where it is white noise; the prior
variance of the coefficient on lag $\ell$ of variable $j$ is proportional to
$\lambda_i^{2}/(\ell^{2}\psi_j)$, with the scales $\psi_j$ fixed at univariate
autoregressive residual-variance estimates. The sum-of-coefficients dummy observations of
\citet{DoanLittermanSims1984} discipline long-run drift with tightness $\mu_i$; on seasonal
blocs the dummy-observation entries for the $D_t$ columns are zeroed and the corresponding
prior rows left diffuse, so the seasonal pattern is estimated essentially from the data.
The two tightnesses are estimated, not calibrated: $\lambda_i$ carries a Gamma hyperprior
with mode $0.2$ and standard deviation $0.4$, $\mu_i$ mode $1$ and standard deviation $1$
--- the settings of \citet{Giannone2015} --- and both are updated by random-walk
Metropolis--Hastings on the marginal likelihood of the completed data, with proposals
reflected into $[10^{-4},5]^2$. The single-unit-root dummy is not used.

The full system is estimated by the following iteration.
\begin{enumerate}\setlength\itemsep{2pt}
  \item \emph{Initialize.} Build starting stars from the raw partner panels (leading gaps
    back-filled with first valid values); set $(\lambda_i,\mu_i)$ at their
    marginal-likelihood modes.
  \item \emph{Bloc pass (data augmentation).} For each bloc $i$, holding its stars fixed:
    alternate (a) drawing the missing cells with the simulation smoother given
    $(\beta_i,\Sigma_i)$; (b) drawing $(\beta_i,\Sigma_i)$ from the normal--inverse-Wishart
    posterior given the completed data; (c) updating $(\lambda_i,\mu_i)$ by
    Metropolis--Hastings. Retain the posterior-mean completed panel.\footnote{Within the
    outer loop the smoother's companion matrix is held fixed at the previous pass's
    posterior mean and common random numbers are used across passes, which makes the outer
    iteration contract in practice. An exact data-augmentation pass, with the companion
    rebuilt at every sweep, is available as an optional final polish.}
  \item \emph{Star rebuild (Jacobi).} Recompute every bloc's foreign variables
    \eqref{eq:rowdem}--\eqref{eq:rowfin} from the partners' just-completed panels.
  \item \emph{Iterate} steps 2--3 until the largest change in any bloc's star history falls
    below $0.03$ (excluding the pandemic-imputation window), with a cap of 12 outer
    iterations.
\end{enumerate}
Each bloc's sampler runs 1{,}000 draws with the first 500 discarded; the shipped model is
the posterior mean $(\hat\beta_i,\hat\Sigma_i)$ together with the 500 retained draws, from
which the tool simulates its uncertainty bands. The scheme is an EM-type algorithm in the
tradition of \citet{DempsterLairdRubin1977} with a Jacobi outer loop across blocs. When it
stops, every completed history, every filled gap, and every foreign variable is a fixed
point of the whole system --- the estimation-time counterpart of the classical model's weak
exogeneity.

\subsection{The solved system}\label{sec:solve}

At solve time the question changes from ``what are the parameters'' to ``what does the
world look like under my conditions.'' A scenario is a set of pinned future cells ---
values a user imposes on particular variables in particular quarters --- and the answer is
the conditional forecast: the smoother of \eqref{eq:ss} run over the future, treating pins
as observations \citep{WaggonerZha1999,BanburaGiannoneLenza2015}. The near-degenerate
measurement variance makes a pin hard; giving a cell a finite variance instead makes it
soft, which is how outside forecasts enter at a chosen level of trust. For a single bloc,
that is the whole story.

Coupled, one more requirement appears: bloc $i$'s foreign variables must equal the
trade-weighted average of what its partners actually do, and the partners' outcomes depend
on their own foreign variables in turn. Because the smoother's conditional mean is $\emph{affine}$
in the imposed values, this consistency requirement is a linear fixed point. Stack into $x$
the deviations from baseline of every bloc's three star paths over the $H=\HorizonQ$-quarter
horizon, augmented by the four global-price paths that the US bloc originates, so that
$\dim x = N\cdot3\cdot H + 4H = 2{,}360$. Let $d$ collect the direct impact of the user's
pins on those objects, and let $M$ be the cross-response operator: entry by entry, how one
bloc's star deviation moves when a partner's solved paths move, built from the partners'
conditional-forecast operators and the weights \eqref{eq:wtilde}. Consistency is
\begin{equation}\label{eq:fixedpoint}
  x \;=\; M x + d
  \qquad\Longrightarrow\qquad
  x \;=\; (I-M)^{-1} d,
\end{equation}
with a unique solution if and only if $\rho(M)<1$. The shipped system has
$\rho(M)=\RhoM$, comfortably inside the unit circle: the margin $1-\rho(M)$ is \RhoMargin{}
percentage points, and it would take a uniform \RhoBreakPct{} percent strengthening of the
entire coupling to reach $\rho=1$. Nearness to one still warrants care, because the worst-case
geometric amplification $1/(1-\rho)$ is \WorldAmpBound{} and steepens sharply toward the
boundary --- a re-estimation lifting $\rho$ to $0.995$ would more than double it --- so
$\rho(M)$ is recomputed and checked at every vintage, and a non-contracting one would not be
shipped; the amplification actually realized in Section~\ref{sec:example} is far smaller, the
pinned shock decaying and the world aggregation reweighting it. The inverse is computed once per
estimation vintage and shipped with the
tool, so at run time a scenario costs one matrix--vector product plus one conditional
forecast per bloc --- which is why the browser answers in about a second. Iterating
\eqref{eq:fixedpoint} bloc by bloc (Gauss--Seidel) converges to the same point and serves
as the verification path.\footnote{Slowly, at this $\rho$: the residual shrinks like
$\rho^{k}$, so hundreds of sweeps are needed --- which is exactly why the inverse is
precomputed rather than iterated live.}

The comparison with Section~\ref{sec:classical} now takes one sentence per object
(Table~\ref{tab:gvar-vs}, Figure~\ref{fig:gbvar}). The classical GVAR stacks the blocs
\emph{contemporaneously} and inverts $G$ once to obtain the world reduced form; the MFFM
stacks the blocs \emph{across the forecast} and inverts $I-M$ once to obtain the world
scenario operator, which is the object a scenario tool needs. Weak exogeneity is replaced
twice over, by the estimation-time fixed point in the star data and the solve-time pinning
of the star paths. The common balanced window is replaced by model-based completion inside
the estimation. Country-by-country reduced-rank regression under rank restrictions is
replaced by hierarchical Bayesian shrinkage in levels. And the US asymmetry is removed: the
United States carries foreign variables like every other bloc, and its four financial
prices are broadcast to all blocs from its solved paths. The \NMembers{} euro-area member
models sit outside this fixed point as satellites: each is a standalone BVAR whose
euro-area columns (the common policy rate, area-wide demand and prices) are pinned to the
EA bloc's solved paths, so members inherit the area outcome without enlarging the coupled
system.

\begin{figure}[t]
\centering
\includegraphics[width=0.92\textwidth]{figs/tikz/fig_gbvar.pdf}
\caption{The MFFM pipeline. The left box is the estimation loop of
Section~\ref{sec:estimation}: complete the ragged panels with the simulation smoother,
re-estimate under the hierarchical prior, rebuild the stars from the completed panels, and
iterate to a fixed point. The United States sits inside the loop as an ordinary coupled
bloc: it originates the global prices and receives the world through its own stars. The
converged companions then feed the affine solve of \eqref{eq:fixedpoint}, whose inverse is
computed once and shipped with the tool.}
\label{fig:gbvar}
\end{figure}

\subsection{What the design buys, and what it costs}\label{sec:tradeoffs}

Why must the system be solved this way? Part of the answer is forced by the data, part is
chosen for the tool, and part is a price paid knowingly. All three deserve to be on the
table.

The data force the estimation design. At \NBlocs{} blocs there is no common window worth
estimating on: samples differ by more than a decade, so truncating every bloc to the shared
span would discard most of the information in the long samples to accommodate the short
ones --- and the pandemic sits in the middle of whatever window remains. A classical
VARX$^{*}$ asks for more observations than a large bloc has in the overlap. Completing the
panel inside the estimation, under shrinkage strong enough to make fifteen-year samples
usable, is what makes a \NBlocs{}-bloc system estimable at all. The EM iteration is then a
requirement rather than a refinement: the stars each bloc needs are functions of the other
blocs' completed panels, so there is nothing to plug in --- consistency between completion
and estimation can only be reached by iterating to it.

The tool chooses the solution design. A scenario instrument must answer \emph{conditional}
questions in seconds, and the classical reduced form answers unconditional ones: $G^{-1}$
delivers dynamics, not the response of a fifty-economy system to ``hold this yield fixed
while that one rises.'' The conditional-forecast operator is the object the tool needs, it
is affine in what the user imposes, and affinity is what turns the coupled solve into
$(I-M)^{-1}d$, computable once. The stars are pinned by deviation at solve time for the
same reason: pinning preserves the linear structure on which interactivity stands.

The gains follow directly. Every economy is estimated on its full sample; the pandemic
enters as missingness rather than distortion; posterior draws deliver bands, not just point
paths; the United States sits inside the system, so the world feeds back; and between
estimations the tool is updated by data operations --- append a quarter, re-anchor the
forecast --- with the model untouched.

The costs deserve equal daylight. The long-run discipline of cointegrating restrictions is
gone, replaced by priors and by user-set long-run anchors; the model will not discover a
theory-consistent equilibrium on its own. Weak exogeneity is not tested but replaced by
construction, so cross-country consistency is imposed rather than examined. The coupled
multiplier stands on $\rho(M)=\RhoM$, close to one, which makes the global amplification
the solve produces powerful and places weight on the star equations that generate it;
Appendix~B confronts the resulting responses with the impulse responses of a classical GVAR
estimated on the same data. And the reduced-form layer identifies nothing by itself:
without the structural overlay of Section~\ref{sec:counterfactual}, a scenario is a
most-likely world, not a statement about cause and effect. The tool is built so that each
of these costs is visible where it bites --- anchors are set by the user, provenance is
flagged on every chart, and the structural layer is drawn as its own line.

\begin{table}[t]
\centering
\caption{The classical GVAR and the MFFM, side by side.}
\label{tab:gvar-vs}
\small
\begin{tabular}{p{0.16\textwidth}p{0.38\textwidth}p{0.38\textwidth}}
\toprule
 & \textbf{Classical GVAR} & \textbf{MFFM} \\
\midrule
Country model &
VARX$^{*}$/VECMX$^{*}$ \eqref{eq:varx}: contemporaneous foreign terms $\Lambda_{i0}$,
cointegrating ranks, long-run restrictions. &
Closed VAR \eqref{eq:mffmbloc} in levels: stars endogenous, co-movement through
$\Sigma_i$, shrinkage instead of rank restrictions. \\
\addlinespace
Estimation &
Each country separately: reduced-rank regression conditional on weakly exogenous stars;
least squares for the short run. &
Hierarchical Bayesian per bloc \citep{Giannone2015}; $(\lambda_i,\mu_i)$ estimated by MH
on the completed-data marginal likelihood. \\
\addlinespace
Foreign variables &
Built once from observed data; fixed weights; weak exogeneity assumed and tested. &
Endogenous columns; rebuilt from partners' completed panels inside the estimation loop;
per-quarter renormalized weights \eqref{eq:wtilde}. \\
\addlinespace
Missing data &
Common balanced window; ragged edges truncated; pandemic handled with dummies. &
Completed inside estimation (EM + simulation smoother); pandemic quarters treated as
missing, financials kept observed. \\
\addlinespace
United States &
Source of $d_t$; restricted foreign-variable set; weak exogeneity fails for the dominant
economy. &
Fully endogenous: carries its own stars; the four global prices broadcast from its solved
paths. \\
\addlinespace
Global solution &
Stack contemporaneously, invert $G$ once \eqref{eq:stack}: the world reduced form. &
Stack across the forecast, invert $I-M$ once \eqref{eq:fixedpoint}: the world scenario
operator; $\rho(M)=\RhoM$. \\
\bottomrule
\end{tabular}
\end{table}

\subsection{The US bloc, concretely}\label{sec:usbloc}

It helps to see all of this in one bloc, and the natural choice is the United States. The
US model is \eqref{eq:mffmbloc} with $n=\NUSVars{}$ variables, $\Gamma_i=0$ (US data are
seasonally adjusted at source), and an estimation sample from \USSampleStart{} to
\USSampleEnd{} --- \USSampleT{} quarters. Its core follows the large US BVAR built at the
Federal Reserve Bank of New York by \citet{CrumpEusepiGiannoneQianSbordone2025}: real
activity and its expenditure components, the labor market, prices, a term structure of
interest rates, corporate bond yields, equity prices and their volatility, housing, and
sentiment; their paper documents each series in detail. To this core the MFFM adds what its
global role requires: the four fiscal drivers, the four global prices (which for the United
States are simply domestic columns), and the three foreign variables. The full roster
appears in Appendix~C.

Partitioning $B_{i\ell}$ and $\Sigma_i$ conformably with \eqref{eq:partition} makes the
transmission channels explicit, and each does something different. The blocks
$B^{x x^{*}}_{i\ell}$ put the three stars, at each of three lags, into every domestic
equation, so a slowdown in partner demand works its way into next quarter's US output,
investment, and employment the way any predictor would. The covariance block
$\Sigma_{xx^{*}}$ correlates foreign and domestic innovations, so a foreign surprise moves
US variables \emph{within} the quarter; when a scenario pins a star path, the conditioning
machinery exploits exactly this. And the stars' own equations, the bottom rows of each
$B_{i\ell}$, close the bloc, replacing the classical weak-exogeneity assumption and letting
the same model run in two modes.

Those two modes are the practical payoff of the design. Used alone, the way a desk
economist uses a country model, the US bloc forecasts its own foreign variables; in effect
it carries a reduced-form view of the world inside itself. In a global solve, its foreign
variables are instead pinned through \eqref{eq:fixedpoint} to what the partner blocs
actually produce, the partners' foreign variables see the United States in turn, and the
fixed point reconciles everyone at once. The difference between the two modes is precisely
the cross-border feedback, and it is measurable: Section~\ref{sec:globalsolve} runs the
credit-spread scenario both ways.

A last word on the fiscal side, because it works differently from the rest of the bloc and
users notice. The VAR forecasts four fiscal \emph{drivers} --- revenue and primary
expenditure as shares of GDP, the effective interest rate on the debt stock, and the
stock-flow residual --- while the debt-to-GDP ratio is not estimated at all. It is computed
quarter by quarter from the government budget identity,
\begin{equation}\label{eq:debt}
  d_t \;=\; d_{t-1}\,
  \frac{1+i_t/400}{1+g_t/100}
  \;-\; \frac{pb_t}{4} \;+\; \frac{s_t}{4},
\end{equation}
where $d_t$ is debt over annualized GDP, $i_t$ the effective interest rate (percent per
annum), $g_t$ nominal GDP growth (percent, quarter-on-quarter, \emph{not} annualized), $pb_t$ the
primary balance in percent of GDP, and $s_t$ the stock-flow adjustment. Every scenario
therefore carries its fiscal arithmetic with it: in a recession scenario revenue falls,
growth denominators shrink, and the debt path steepens because accounting says it must, not
because a debt equation was fitted. Section~\ref{sec:example} shows this at work, and the
same identity run backwards lets a user impose a debt target and read off the primary
balance that would deliver it.
